\section*{The Configuration Space of Knots}

To derive the mass spectrum we define the complete geometric state space of the $E_8$ lattice projection, the set of all mathematically valid lattice excitations and then apply physical constraints to determine which configurations can persist. From this list of persistent excitations we will check if we listed known particles, hadrons, etc.

The configuration space is defined by a \textbf{Harmonic Series} (the resonant tiers) modified by \textbf{Geometric Operators} (the projection costs).

\textit{Note on Scope:} This configuration space generates \textbf{Topological Knots} (Fermions/Bosons). \textbf{Lattice Phonons} (Neutrinos) follow the Inverse Resonance hierarchy ($M \propto \Delta^{-n}$) described in Section VIII and are excluded from this specific spectral scan.

\subsection*{The Harmonic Tiers (\texorpdfstring{$\Delta^n$}{Deltapown})}

The fundamental frequency of the lattice is the Heegner Resonance $\Delta = 43$. Just as a vibrating string supports integer harmonics, the lattice supports \textbf{Geometric Powers} of this resonance. These powers define the \textbf{Tier} (or ``Order of Magnitude'')of the excitation.

We identify four distinct resonant tiers:

\begin{table}[h]
\centering
\begin{tabular}{llll}
\toprule
\textbf{Tier} & \textbf{$m_e$} & \textbf{Geometry} & \textbf{Candidate Domain} \\ \midrule
\textbf{0} & $\Delta^0 = 1$ & Unitary Monopole & Electron \\
\textbf{1} & $\Delta^1 = 43$ & Vector (String) & Muon, Pions \\
\textbf{2} & $\Delta^2 = 1,849$ & Tensor (Membrane) & Proton, Tau, Charm \\
\textbf{3} & $\Delta^3 = 79,507$ & Volumetric (Bulk) & Top, Weak Bosons \\ \bottomrule
\end{tabular}
\end{table}

\textit{Note on Units:} The fundamental mass unit of the lattice ($m_e$) is not an arbitrary input. Using the geometric resolution floor derived in Paper I ($m_e \approx PM_{geo}$), the electron represents the unitary ground state against which all other resonant tiers are scaled.

\subsection*{The Geometric Operators}

To generate the candidate spectrum, we define the allowed geometric operations on the fundamental resonance $\Delta^n$. The Configuration Space $\mathbb{S}$ is constructed from three classes of operators derived from the lattice interaction logic:

\textbf{1. Linear Modifiers (The Additive Terms)} \\
These terms represent the static cost of embedding a topology into the manifold.
\begin{itemize}
    \item \textbf{Manifold Projection ($\pi D, \pi$):} The cost of confining flux within 4 spacetime dimensions ($D$) or defining a circular boundary ($\pi$).
    \item \textbf{Lattice Invariants ($\sigma, \chi, \nu$):} The energy required to maintain the lattice symmetry ($\sigma=5$), topological lock ($\chi=2$), or chiral capacity ($\nu=16$).
    \item \textbf{Scalar Anchors ($D, \sigma$):} In Tier 0 (Point Geometry), the manifold dimension ($D=4$) and symmetry order ($\sigma=5$) can act as fundamental mass units, representing the geometric cost of anchoring a point to the spacetime structure.
\end{itemize}

\textbf{2. The Saturation Scalars (The Multipliers)} \\
While fundamental excitations scale as integer powers, resonances that approach the saturation limit of their dimensional container require a \textbf{Projection Scalar ($S$)} to quantify the manifold overflow.
\begin{itemize}
    \item \textbf{Areal Saturation ($S_b$):} Applies to Tier 2. Represents the membrane filling the spacetime dimension $D$ and overflowing into the symmetry sector $\chi/\sigma$.
    \[ S_b = D + \frac{\chi}{\sigma} = 4 + 0.4 = \mathbf{4.4} \]
    \textit{(This scalar generates the Bottom Quark base scale: $1849 \times 4.4 \approx 8135 m_e$)}.
    \item \textbf{Volumetric Saturation ($S_t$):} Applies to Tier 3. Represents the volume filling the manifold $D$ and overflowing reciprocal leakage $1/D$.
    \[ S_t = D + \frac{1}{D} = 4 + 0.25 = \mathbf{4.25} \]
    \textit{(This scalar generates the Top Quark base scale: $79507 \times 4.25 \approx 337,900 m_e$)}.
\end{itemize}

\textbf{3. The Symmetry Half-Step (Isospin Projection)} \\
Generally, geometric modifiers utilize integer coefficients $k \in \mathbb{Z}$ representing full rotations of the lattice symmetry. However, the Selection Principle permits a specific relaxation for topological doublets (Isospin).
\begin{itemize}
    \item \textbf{Mechanism:} The proton and neutron form an $SU(2)$ isospin doublet with components $I_3 = \pm 1/2$. When this internal quantum number is projected onto the 5-fold symmetry of the lattice ($\sigma=5$), it generates a geometric mass shift proportional to the product of the spin projection and the symmetry order:
    \[ \Delta M = I_3 \times \sigma = \pm \frac{1}{2} \times 5 = \pm 2.5 \, m_e \]
    \item \textbf{Physical Constraint:} This half-integer rotation creates a \textbf{Geometric Phase Offset} that decouples the knot from the electromagnetic boundary constraint ($\chi$). Consequently, the fractional coefficient \textbf{$k = \pm 0.5$} is valid if and only if the resulting state is \textbf{electrically neutral} (e.g., the Neutron). This explains why the Proton cannot access this lower-energy geometry.
\end{itemize}

\textbf{4. Resonance Multiplicity ($c$)} \\
The Fundamental Harmonic $\Delta^n$ represents a single standing wave in the lattice. However, the 5-fold symmetry of the substrate ($\sigma=5$) allows for \textbf{Tier Superposition}, where the lattice supports integer multiples of the base frequency up to the symmetry limit.
\begin{itemize}
    \item \textbf{Definition:} $c \in \{1, 2, \dots, \sigma\}$.
    \item \textbf{Physical Meaning:} This represents the \textbf{Symmetry Occupancy} of the state—how many distinct symmetry sectors the resonance occupies simultaneously.
    \item \textbf{Example:} The Tau Lepton ($M \approx 3477 m_e$) is identified as a \textbf{Double Occupancy} excitation of the Areal Tier ($c=2, n=2$), creating a base resonance of $2\Delta^2 = 3698 m_e$ before drag corrections.
\end{itemize}

\textbf{5. Composite Modifiers (Interaction Products \& Impedance)} \\
Beyond simple linear dimensions, the lattice supports composite operators representing the interaction of distinct geometric sectors.
\begin{itemize}
    \item \textbf{Symmetry Drag ($\sigma \Delta$):} The friction of a linear string ($\Delta$) rotating against the 5-fold symmetry ($\sigma$). This product ($\approx 215 m_e$) is the dominant mass term for the Muon.
    \item \textbf{Electroweak Saturation ($\nu \sigma$):} The product of Channel Capacity ($\nu$) and Interaction Order ($\sigma$). This term ($\nu\sigma = 80$) represents the baseline load of the electroweak coordination sector.
    \item \textbf{The Vacuum Impedance ($\alpha^{-1}$):} While Fermions are resonances of the lattice geometry, Mesons (Pions, Kaons) function as excitations of the binding field. The configuration space includes the \textbf{Aggregate Vacuum Impedance} ($\alpha^{-1} \approx 137.036$). It serves as the Base Tier for light mesons (e.g., $m_{\pi} \approx 2\alpha^{-1}$).
\end{itemize}

\subsection*{The Generative Equation}

We formally define the \textbf{Universal Mass Function} used for the blind spectral analysis. Any candidate mass $M$ must satisfy:

\[ M = \left( c \cdot S \cdot \Delta^n \right) + \sum_{i=1}^{N_{max}} k_i \mathcal{O}_i \]

Where:
\begin{itemize}
    \item \textbf{Search Domain:} To enforce low entropy, we restrict Complexity to $N_{max} \le 3$ (terms) and Coefficients to $k_i \in \{-5, \dots, 5\}$. \textbf{Exception:} For the Symmetry Operator ($\sigma$), include the fractional coefficient \textbf{$k = \pm 0.5$}. This represents the Isospin projection ($I_z = \pm 1/2$) acting on the 5-fold symmetry, yielding a mass shift of $\pm 2.5$.
    \item \textbf{Parameters:}
    \begin{itemize}
        \item $\Delta^n$: Tier $n \in \{0, 1, 2, 3\}$.
        \item $c$: Multiplicity $\{1, 2, 3, 4, 5\}$.
        \item $S$: Scalar constraint ($S=1$ for $n<2$; $S \in \{1, S_b\}$ for $n=2$; $S \in \{1, S_t\}$ for $n=3$).
    \end{itemize}
\end{itemize}